\documentclass[a5paper,10pt]{article}

% https://github.com/InDevRus/filippov-solutions

% Нумерация страниц
\pagenumbering{gobble}

% Настройка полей
\usepackage{geometry}
\geometry{left=1cm}
\geometry{right=1cm}
\geometry{top=1cm}
\geometry{bottom=1.5cm}

% Вставка таблицы
\usepackage{array}
\usepackage{tabularx}

% Инструменты выравнивания формул
\usepackage{amsmath}
\usepackage{mathtools}

% Диакритика в математике
\usepackage{amsfonts}

% Вычисления размеров на ходу
\usepackage{calc}

% Удобные записи производных
\usepackage[italic=true]{derivative}

% Настройки сносок
\usepackage{enotez}

% Длинные знаки векторов
\usepackage{anyfontsize}
\usepackage[b]{esvect}

% Вставка картинок
\usepackage{graphicx}

% Вставка красной строки
\usepackage{indentfirst}
\setlength{\parindent}{1em}

% Условия в командах
\usepackage{xifthen}

% Луа-скрипты
\usepackage{luacode}

% Настройка команд отдельно в математическом и текстовом режимах
\usepackage{mathcommand}

% Для повторения LaTeX кода
\usepackage{multido}

% Конфигурация отступов формул
\delimitershortfall-1sp
\usepackage{mleftright}
\mleftright

% Растягивание объектов
\usepackage{relsize}

% Обтекание текстом
\usepackage{wrapfig}
\usepackage{subcaption}
\usepackage{floatflt}

% Поддержка ввода кириллицы
\usepackage[english, russian]{babel}

% Настройка корневой позиции для компилятора
\usepackage{import}
\providecommand*{\ProjectRootPath}{..}

\import{\ProjectRootPath/preambles/macros/}{theorems}

\import{\ProjectRootPath/preambles/fonts}{font_setup}

% Знак градуса
\usepackage{siunitx}

\import{\ProjectRootPath/preambles/macros/}{operators}
\import{\ProjectRootPath/preambles/macros/}{redefinitions}

\import{\ProjectRootPath/preambles/fonts}{letters_setup}
\import{\ProjectRootPath/preambles/fonts/adjustments}{custom_superscripts}

\import{\ProjectRootPath/preambles/custom}{formatting}
\import{\ProjectRootPath/preambles/custom}{frames}
\import{\ProjectRootPath/preambles/custom}{list_setup}

% TODO: перевести команды на современный стиль объявления
\input{../preambles/plotting}

\begin{document}

$\pow{e}{-s}\br{1 + \der{s}{t}} = 1$, где $s = s\br{t}$.
$\der{s}{t} = \pow{e}{s} - 1$.
$\dfrac {s\`\br{t}} {\pow{e}{s} - 1} = 1$,
а функция $s\br{t} = 0$ относится к решениям.
\begin{align*}
    \tint \dfrac {s\`\br{t}} {\pow{e}{s} - 1} \di t
    = \tint \dfrac {\di s} {\pow{e}{s} - 1} =
    \begin{bmatrix}
        s = \ln u \\
        \di s = \dfrac {du} {u}
    \end{bmatrix} 
    = \tint \dfrac {du} {u\br{u - 1}}
    & = \tint \br{\dfrac {1} {u - 1} - \dfrac {1} {u}} du \\
    & = \ln\vbr{u - 1} - \ln\vbr{u} + C \\
    & = \ln\vbr{1 - \pow {e}{-s}} + C.
\end{align*}

Итак, взятие интеграла обеих частей уравнения 
$\dfrac {s\`\br{t}} {\pow{e}{s} - 1} = 1$.
даёт $\ln\vbr{\pow {e}{-s} - 1} = t + C$.
$\pow{e}{-s} = 1 + C\pow{e}{t}$.
$s\br{t} = -\ln\vbr{1 + C\pow{e}{t}}$,
а функция $s\br{t} = 0$ входит в такое общее решение при $C = 0$.

\begin{figure}[!h]
\begin{minipage}{0.5\textwidth}
    \centering
    \begin{tikzpicture}
        \pgfplotsset{
            Graph/.style = {
                samples = 200,
            },
        }
        \begin{axis}[
            width = 7.5cm,
            xlabel = $t$,
            ylabel = $s$,
            ymin = -4.2,
            ymax = 4.2,
            xmin = -5.2,
            xmax = 5.2, 
            xtick = {-10, ..., 10},
            ytick = {-10, ..., 10},
            minor tick num = 3,
            declare function = {
                f(\C, \x) = -ln(abs(1 + \C * exp(x)));
            }]
            
            \draw[thick, blue] (-5.5, 0) -- (5.5, 0);
            
            \pgfplotsinvokeforeach {0.01, 0.05, 0.1, 0.25, 0.5, 1, 3, 10, 50} {
                \addplot[Graph, domain = -5.5 : 5.5]
                {f(#1, x)};
            }
        \end{axis}
    \end{tikzpicture}
    \caption*{$C > 0$}
\end{minipage}
\begin{minipage}{0.5\textwidth}
    \centering
    \begin{tikzpicture}
        \pgfplotsset{
            Graph/.style = {
                samples = 200,
            },
        }
        \begin{axis}[
            width = 7.5cm,
            xlabel = $t$,
            ylabel = $s$,
            ymin = -4.2,
            ymax = 4.2,
            xmin = -5.2,
            xmax = 5.2,
            xtick = {-10, ..., 10},
            ytick = {-10, ..., 10},
            minor tick num = 3,
            declare function = {
                f(\C, \x) = -ln(abs(1 + \C * exp(x)));
            }]

            \draw[thick, blue] (-5.5, 0) -- (5.5, 0);

            \pgfplotsinvokeforeach {-50, -10, -2, -1, -0.5, -0.25, -0.05} {
                \addplot[Graph, domain = -5.5 : -ln(abs(#1)) - 0.01]
                {f(#1, x)};
                \addplot[Graph, domain = -ln(abs(#1)) + 0.01 : 5.5]
                {f(#1, x)};
            }
        \end{axis}
    \end{tikzpicture}
    \caption*{$C < 0$}
\end{minipage}
\end{figure}

\end{document}
